\section*{Hoofdstuk \ref{ch:b2:mantle-melting} --- De mantel smelten}

\begin{solution}{exo:b2:mantle-melting:1}
(a) Meer dan \qty{90}{\percent} olivijn: \omterm{def:b2:mantle-melting:peridotites}{duniet}. (b) Olivijn en orthopyroxeen
zonder clinopyroxeen: \omterm{def:b2:mantle-melting:peridotites}{harzburgiet}. (c) Meer dan \qty{40}{\percent} olivijn met
beide pyroxenen: \omterm{def:b2:mantle-melting:peridotites}{lherzoliet}. (d) Minder dan \qty{40}{\percent} olivijn: een pyroxeniet
(een olivijnwebsteriet, met beide pyroxenen). (e) $57/98 = \qty{58.2}{\percent}$,
$28/98 = \qty{28.6}{\percent}$, $13/98 = \qty{13.3}{\percent}$: een \omterm{def:b2:mantle-melting:peridotites}{lherzoliet}.
\end{solution}

\begin{solution}{exo:b2:mantle-melting:2}
(a) De liquidus bij \qty{60}{\percent} anorthiet ligt op \qty{1390}{\celsius}, in het
anorthietveld: anorthiet kristalliseert als eerste. (b) Diopsied verschijnt in het
eutecticum, \qty{1270}{\celsius}, wanneer de smelt \qty{42}{\percent} anorthiet bevat.
(c) De eerste smelt van elk mengsel van de twee mineralen is de eutectische smelt,
\qty{42}{\percent} anorthiet, bij \qty{1270}{\celsius}.
\end{solution}

\begin{solution}{exo:b2:mantle-melting:3}
(a) $\num{4.0e-5}\times1623/(3300\times1000) = \qty{1.97e-8}{K/Pa} =
\qty{19.7}{K/GPa}$. (b) $\num{e9}/(3300\times9.81) = \qty{3.09e4}{m}$, of
\qty{30.9}{km} per gigapascal; $19.7/30.9 = \qty{0.64}{K/km}$, ongeveer vijftig keer
minder dan de \qty{30}{K/km} dicht bij het oppervlak, waar warmte wordt geleid, niet
meegevoerd.
\end{solution}

\begin{solution}{exo:b2:mantle-melting:4}
(a) $T_{\mathrm s}(0) = \qty{1085.7}{\celsius}$; $T_{\mathrm s}(1) = 1085.7 +
132.9 - 5.1 = \qty{1213.5}{\celsius}$; $T_{\mathrm s}(2) = 1085.7 + 265.8 - 20.4 =
\qty{1331.1}{\celsius}$; $T_{\mathrm s}(3) = 1085.7 + 398.7 - 45.9 =
\qty{1438.5}{\celsius}$. (b) $132.9 - 2\times5.1\times2 = \qty{112.5}{K/GPa}$.
(c) Bij \qty{3}{GPa} ligt \qty{1400}{\celsius} onder de solidus,
\qty{1438.5}{\celsius}; bij het opstijgen tot \qty{2}{GPa} langs de adiabaat koelt de mantel maar
ongeveer \qty{20}{\celsius} af, tot zo'n \qty{1380}{\celsius}, terwijl de solidus daalt tot
\qty{1331.1}{\celsius}: het gesteente ligt nu boven zijn solidus en een deel ervan smelt.
\end{solution}

\begin{solution}{exo:b2:mantle-melting:5}
(a) $1350 + (1277.5 - 1350)\times(30 - 20)/(40 - 20) = \qty{1314}{\celsius}$.
(b) De samenstelling ligt aan de diopsiedkant van het eutecticum: diopsied.
(c) De smelt bereikt het eutecticum met \qty{42}{\percent} anorthiet; anorthiet
blijft in de smelt, dus de \omterm{def:b2:mantle-melting:fraction}{smeltfractie} is $30/42 = 0.714$ en $1 - 0.714 =
\qty{29}{\percent}$ van het systeem is als diopsied gekristalliseerd. (d) Alles
kristalliseert: \qty{70}{\percent} diopsied en \qty{30}{\percent} anorthiet, de
totale samenstelling.
\end{solution}

\begin{solution}{exo:b2:mantle-melting:6}
(a) Forsteriet: $60.08/(2\times40.30 + 60.08) = 60.08/140.68 =
\qty{42.71}{\percent}$ \ce{SiO2}; enstatiet: $60.08/100.38 = \qty{59.85}{\percent}$.
(b) $0.975\times59.85 + 2.5 = \qty{60.85}{\percent}$. (c) Enstatiet ligt tussen
forsteriet en de smelt: de fractie forsteriet is $(60.85 - 59.85)/(60.85 -
42.71) = 1.00/18.14 = \qty{5.5}{\percent}$, de waarde van Bowen en Andersen. (d) De
smelt is rijker aan silica dan het orthopyroxeen waaruit hij komt, en het residu
wint olivijn.
\end{solution}

\begin{solution}{exo:b2:mantle-melting:7}
(a) $5.1P^2 - 112.9P + 214.3 = 0$: $P_0 = (112.9 - \sqrt{12746.4 -
4371.7})/10.2 = (112.9 - 91.51)/10.2 = \qty{2.10}{GPa}$, op $2.097\times30.9 =
\qty{64.8}{km}$. (b) $5.1P^2 - 112.9P + 364.3 = 0$: $P_0 = (112.9 -
72.90)/10.2 = \qty{3.92}{GPa}$, op \qty{121}{km}. (c) $T_{\mathrm s}'(2.10) = 132.9
- 21.4 = \qty{111.5}{K/GPa}$, dus $\Delta P_0 = 10/(111.5 - 20) = \qty{0.109}{GPa}$
per \qty{10}{\celsius}, of \qty{3.4}{km}.
\end{solution}

\begin{solution}{exo:b2:mantle-melting:8}
(a) $T_{\mathrm s} = \qty{1331.1}{\celsius}$, $T_{\mathrm l} = 1475.0 + 160.0 -
12.8 = \qty{1622.2}{\celsius}$; $T = 1331.1 + 0.10\times291.1 =
\qty{1360.2}{\celsius}$, of \qty{1633}{K}. (b) $(\partial T/\partial P)_F = 112.5
+ 0.10\times(67.2 - 112.5) = \qty{108.0}{K/GPa}$; $(\partial T/\partial F)_P =
T_{\mathrm l} - T_{\mathrm s} = \qty{291.1}{K}$. (c) $c_p/T = 0.612$; met $\alpha/\rho
= \qty{12.1}{J/(kg.K)}$ per gigapascal is de teller $0.612\times108.0 - 12.1
= 54.0$ en de noemer $300 + 0.612\times291.1 = 478.2$: $-\mathrm dF/\mathrm
dP = 0.113$, \qty{11.3}{\percent} per gigapascal. (d) $54.0/178.2 = 0.303$,
\qty{30}{\percent} per gigapascal: de latente warmte verlaagt de productiviteit met een
factor 2.7.
\end{solution}

\begin{solution}{exo:b2:mantle-melting:9}
(a) $43\times1 = \qty{43}{\celsius}$; $43\times5^{0.75} = 43\times3.34 =
\qty{144}{\celsius}$; $43\times13^{0.75} = 43\times6.85 = \qty{294}{\celsius}$.
(b) $1213.5 - 294 = \qty{919}{\celsius}$. (c) \qty{125}{ppm} is
\qty{0.0125}{\percent}, dus $X = 0.0125/0.01 = \qty{1.25}{\percent}$ en $\Delta T =
43\times1.25^{0.75} = 43\times1.18 = \qty{51}{\celsius}$.
\end{solution}

\begin{solution}{exo:b2:mantle-melting:10}
(a) De basis is $2H$ breed en de mantel doorkruist ze met $U$: $2UH$ per lengte-eenheid
rug en per tijdseenheid. Elke zijde overspant de hoogte $H$ en wordt horizontaal doorkruist met
$U$: $UH$ elk, $2UH$ in totaal. Met zijden onder $\theta$ zou de basis
$2H/\tan\theta$ breed zijn en de instroom $2UH/\tan\theta$, alleen gelijk aan de uitstroom als
$\tan\theta = 1$: in dit model wordt de hoek door massabehoud vastgelegd. (b) Met
constante productiviteit is de \omterm{def:b2:mantle-melting:fraction}{smeltfractie} aan de top $F_{\max} = \Pi P_0 = \Pi\rho
g z_0$, en $H = z_0$: $h_c = F_{\max}z_0/2 = \Pi\rho g z_0^2/2$. (c) Elk jaar
\qty{50}{mm} nieuwe plaat, $\qty{0.05}{m}\times\qty{1000}{m}\times\qty{7100}{m} =
\qty{3.6e5}{m^3}$ korst per kilometer rug; bij een gemiddelde \omterm{def:b2:mantle-melting:fraction}{smeltfractie} van \qty{12}{\percent}
komt die uit $\num{3.6e5}/0.12 = \qty{3.0e6}{m^3}$ mantel.
\end{solution}

\begin{solution}{exo:b2:mantle-melting:11}
(a) Nadat een massa $t$ aan A een eenheidsmassa smelt heeft verlaten, bevat de smelt $0.6 -
t$ aan A, 0.2 aan B en 0.2 aan C in een massa $1 - t$: $\big((0.6 - t), 0.2,
0.2\big)/(1 - t)$, en B/C blijft 1. (b) Op de \omterm{def:b2:mantle-melting:ternary}{cotectische lijn} geldt $(0.40 - 0.15s, 0.60 -
0.30s, 0.45s)$; B = C geeft $0.60 - 0.30s = 0.45s$, $s = 0.8$: $L_1$ heeft A 28, B
36, C 36. Dan $0.2/(1 - t) = 0.36$, $t = 0.44$: \qty{44}{\percent} van het systeem is
als A gekristalliseerd. (c) C zit volledig in de smelt: $f = 20/45 = 0.44$,
\qty{44}{\percent} nog vloeibaar. (d) A: $0.60 - 0.444\times0.25 = 0.49$; B: $0.20 -
0.444\times0.30 = 0.07$; met de smelt erbij $49 + 7 + 44 = \qty{100}{\percent}$.
\end{solution}

\begin{solution}{exo:b2:mantle-melting:12}
(a) $X_{\mathrm{sat}}(1) = 12 + 1 = \qty{13}{\percent}$, $\Delta T =
\qty{294}{\celsius}$, verzadigde solidus $1213.5 - 294 = \qty{919}{\celsius}$;
$X_{\mathrm{sat}}(2) = 12\times1.516 + 2 = \qty{20.2}{\percent}$, $\Delta T =
43\times20.2^{0.75} = 43\times9.53 = \qty{410}{\celsius}$, verzadigde solidus $1331.1
- 410 = \qty{921}{\celsius}$. (b) De smelt lost bij hogere druk meer water op,
en de grotere verlaging heft de stijging van \qty{118}{\celsius} van de droge solidus op.
(c) De solidus moet dalen met $1331.1 - 1100 = \qty{231}{\celsius}$: $X =
(231/43)^{1/0.75} = 5.37^{1.333} = \qty{9.4}{\percent}$ water, minder dan de
\qty{20.2}{\percent} bij verzadiging, dus de wig kan smelten.
\end{solution}

\begin{solution}{pb:b2:mantle-melting:1}
\textbf{1.} $1350 + 273 = \qty{1623}{K}$.
\textbf{2.} $\num{4.0e-5}\times1623/(3300\times1000) = \qty{1.97e-8}{K/Pa} =
\qty{19.7}{K/GPa}$, afgerond tot \qty{20}{K/GPa}.
\textbf{3.} $\num{e9}/(3300\times9.81) = \qty{30.9}{km}$; $20/30.9 =
\qty{0.65}{K/km}$.
\textbf{4.} $1350 + 20P = 1085.7 + 132.9P - 5.1P^2$, dat wil zeggen $5.1P^2 - 112.9P +
264.3 = 0$: $P_0 = (112.9 - \sqrt{12746.4 - 5391.7})/10.2 = (112.9 - 85.76)/10.2 =
\qty{2.66}{GPa}$.
\textbf{5.} $2.66\times30.9 = \qty{82}{km}$; $T = 1350 + 20\times2.66 =
\qty{1403}{\celsius}$.
\textbf{6.} $1085.7 + 132.9\times2.66 - 5.1\times2.66^2 = 1085.7 + 353.5 - 36.1 =
\qty{1403}{\celsius}$.
\textbf{7.} $T_{\mathrm s} = 1085.7 + 199.4 - 11.5 = \qty{1273.6}{\celsius}$;
$T_{\mathrm l} = 1475.0 + 120.0 - 7.2 = \qty{1587.8}{\celsius}$; verschil
\qty{314.2}{K}.
\textbf{8.} $1273.6 + 0.10\times314.2 = \qty{1305.0}{\celsius}$, \qty{1578}{K}.
\textbf{9.} $T_{\mathrm s}' = 132.9 - 15.3 = \qty{117.6}{K/GPa}$, $T_{\mathrm l}' =
80.0 - 9.6 = \qty{70.4}{K/GPa}$: $(\partial T/\partial P)_F = 117.6 +
0.10\times(70.4 - 117.6) = \qty{112.9}{K/GPa}$; $(\partial T/\partial F)_P =
\qty{314.2}{K}$.
\textbf{10.} $c_p/T = 1000/1578 = 0.634$; teller $0.634\times112.9 - 12.1 = 71.6
- 12.1 = 59.5$; noemer $300 + 0.634\times314.2 = 499.2$: $0.119$, of
\qty{11.9}{\percent} per gigapascal.
\textbf{11.} Met de gefitte functie $F = T'^{1.5}$, waarbij $T'$ het deel van de weg van solidus naar liquidus is, moet de temperatuur steil stijgen met $F$ om
de eerste procent smelt te maken, dus $(\partial T/\partial F)_P$ is groot dicht bij de solidus
en het smelten begint traag. $F_{\max} = 0.09\times2.66 = 0.24$, \qty{24}{\percent}.
\textbf{12.} $0.13/0.50 = 0.26$: het clinopyroxeen raakt op bij \qty{26}{\percent}
smelten. Bij \qty{24}{\percent} bevat het residu onder de as er nog een beetje van:
een \omterm{def:b2:mantle-melting:peridotites}{lherzoliet} arm aan clinopyroxeen, op de grens met \omterm{def:b2:mantle-melting:peridotites}{harzburgiet}.
\textbf{13.} $h_c = 0.24\times82/2 = \qty{9.8}{km}$.
\textbf{14.} $F_{\max}/2 = \qty{12}{\percent}$.
\textbf{15.} \qty{9.8}{km} is \qty{38}{\percent} meer dan \qty{7.1}{km}, boven het
bereik van normale korst.
\textbf{16.} Constante productiviteit geeft de diepste mantel te veel smelt; het
gefitte model maakt de eerste smelt traag, $F(z)$ blijft lager dicht bij het begin, en
$\int F\,\mathrm dz$ is kleiner.
\textbf{17.} $F_{\max} = \Pi P_0$ en $H = 30.9P_0$ (km): $h_c = \Pi P_0\times30.9P_0/2
= kP_0^2$ met $k = 0.09\times30.9/2 = \qty{1.39}{km/GPa^2}$.
\textbf{18.} $P_0 = \sqrt{7.1/1.39} = \sqrt{5.11} = \qty{2.26}{GPa}$.
\textbf{19.} $2.26\times30.9 = \qty{70}{km}$.
\textbf{20.} $T_{\mathrm s}(2.26) = 1085.7 + 300.4 - 26.0 = \qty{1360.0}{\celsius}$,
en $T_p = 1360.0 - 20\times2.26 = 1360.0 - 45.2 = \qty{1315}{\celsius}$.
\textbf{21.} Voor \qty{6.3}{km}: $P_0 = \sqrt{4.53} = \qty{2.13}{GPa}$, $T_{\mathrm s}
= \qty{1345.6}{\celsius}$, $T_p = \qty{1303}{\celsius}$. Voor \qty{7.9}{km}: $P_0 =
\sqrt{5.68} = \qty{2.38}{GPa}$, $T_{\mathrm s} = \qty{1373.1}{\celsius}$, $T_p =
\qty{1326}{\celsius}$. De spreiding van de korst laat $1315 \pm
\qty{12}{\celsius}$ toe.
\textbf{22.} $P_0 = \sqrt{20/1.39} = \sqrt{14.39} = \qty{3.79}{GPa}$; $T_{\mathrm s} =
1085.7 + 503.7 - 73.3 = \qty{1516.1}{\celsius}$; $T_p = 1516.1 - 75.8 =
\qty{1440}{\celsius}$, een overschot van $1440 - 1315 = \qty{125}{\celsius}$.
\textbf{23.} De twee gebruiken verschillende gegevens. De lava's leggen de heetste mantel vast,
in het centrum van de pluim; de korst middelt de smelten van alle mantel die
door het smeltgebied gaat, waarvan veel koeler is. Het eenvoudige model houdt ook
de productiviteit constant, hoewel het clinopyroxeen opraakt bij \qty{26}{\percent},
terwijl hier $\Pi P_0 = 0.09\times3.79 = 0.34$: echte hete mantel smelt minder, en heeft
een hogere temperatuur nodig om dezelfde korst te maken.
\textbf{24.} $T_{\mathrm s}'(2.26) = 132.9 - 23.1 = \qty{109.8}{K/GPa}$; $\Delta P_0 =
10/(109.8 - 20) = \qty{0.111}{GPa}$; omdat $h_c \propto P_0^2$, is $\Delta h_c/h_c =
2\times0.111/2.26 = \qty{9.8}{\percent}$, ongeveer een tiende per \qty{10}{\celsius}.
\textbf{25.} Bij een zeer trage rug stijgt de mantel zo langzaam op dat hij warmte verliest door
geleiding naar het koude oppervlak; hij verlaat de adiabaat, zijn smelten stopt
dieper, en de top van de driehoek, waar $F$ het grootst is, wordt afgesneden: de korst is
dunner.
\textbf{26.} De korst van normale ruggen geeft een \omterm{def:b2:mantle-melting:tp}{potentiële manteltemperatuur} van
\textbf{\qty[text-series-to-math=true]{1315}{\celsius}}, binnen ongeveer
\qty{12}{\celsius}: \qty{35}{\celsius} onder de \qty{1350}{\celsius} uit de
mineraalfysica, die met ongeveer evenveel onzeker is, en ruim binnen het bereik 1280 tot
\qty{1400}{\celsius} uit lava's. Drie onafhankelijke thermometers zijn het eens over een mantel
rond \qty{1300}{\celsius} tot \qty{1350}{\celsius}.
\end{solution}
